%%%PAGE 224 rot=0 legibility=good summary=描述运动方法（语言描述法、图象描述：斜率与面积及p-V循环例、公式法：匀变速公式、斜面上滑下滑、中时速与中位速）
%%%BLOCK id=224-01 ch=01 sec=描述运动方法 kind=summary alt=04,08 cont=none y=0.05-0.25
\textbf{描述运动方法}

1）语言描述法
\begin{itemize}
\item 匀速运动
\item 匀变速直线运动
\item 匀变速曲线运动\quad 平抛/斜抛
\item 匀速圆周运动
\item 简谐运动
\end{itemize}
%%%END
%%%BLOCK id=224-02 ch=01 sec=x-t、v-t、a-t图像 kind=knowledge alt=90 cont=none y=0.26-0.35
2）图象描述

斜率 $K=\dfrac{\Delta y}{\Delta x}$\quad 斜率正负不表示大小，表示方向

面积 $S=\int xy$\quad 图象与横轴组成的面积\quad 不表正负，表方向 % CHECK: 原文“不表正负，表方向”，疑漏字（对照上一行“正负不表示大小，表示方向”）
%%%END
%%%BLOCK id=224-03 ch=16 sec=热力学定律·p-V图像 kind=problem alt=90 cont=none y=0.36-0.49
\begin{problem}
%FIG p224_a 0.285 0.36 0.465 0.433
\notefig[0.25]{figures/p224_a.png}
%FIG p224_b 0.478 0.36 0.565 0.412
\notefig[0.25]{figures/p224_b.png}
\begin{solution}
$1\unit{atm}=1\times10^5\unit{Pa}$

$W=$\unclear{…}$(S_1-S_2)\times(-1)$ % 原稿：“=”后有一处被涂划的字迹

$W=-4000\unit{J}$\quad $\Delta U=Q+W=0$\quad 故吸热 $4000\unit{J}$
\end{solution}
\end{problem}
功的正负\quad 在 顺吸 逆放
%%%END
%%%BLOCK id=224-04 ch=01 sec=x-t、v-t、a-t图像 kind=note alt=none cont=none y=0.49-0.52
$v$-$t$、$x$-$t$。
%%%END
%%%BLOCK id=224-05 ch=01 sec=匀变速直线运动公式 kind=formula alt=none cont=none y=0.52-0.63
\textbf{3）}公式法\quad —匀变速直线运动 $V=V_0+at$
%FIG p224_c 0.495 0.505 0.62 0.56
\notefig[0.25]{figures/p224_c.png}
$V_{t/2}=V_0+\dfrac12at$
%FIG p224_d 0.595 0.585 0.745 0.625
\notefig[0.25]{figures/p224_d.png}
%%%END
%%%BLOCK id=224-06 ch=01 sec=匀变速直线运动公式 kind=formula alt=none cont=none y=0.58-0.72
匀变速直线运动\quad $a$ 恒定/恒力

函数关系。系数涵义
\begin{align*}
x&=V_0t+\frac12at^2 &&\text{抛物线}\ \checkmark\\
V&=V_0+at &&\checkmark\ \text{基本\unclear{短}试}\\
V_t^2-V_0^2&=2ax &&\checkmark % CHECK: 首项 V 后似有 t 的笔画，可能为 V^2
\end{align*}
%%%END
%%%BLOCK id=224-07 ch=03 sec=斜面·上滑与下滑 kind=method alt=01 cont=none y=0.72-0.86
%FIG p224_e 0.065 0.745 0.295 0.84
\notefig[0.3]{figures/p224_e.png}
$\mu<\tan\alpha$
\[
\frac{0^2-V_0^2=-2a_1x}{V^2-0^2=2a_2x}\qquad
\left(\frac{V_0}{V}\right)^2=\frac{a_1}{a_2}=\frac{g\sin\theta+\mu\cos\theta}{\sin\theta-\mu\cos\theta}
\] % CHECK: 分子“g sinθ+μcosθ”与分母不对应（g 的位置）；图中倾角记 α，公式中写 θ
%FIG p224_f 0.55 0.815 0.65 0.86
\notefig[0.25]{figures/p224_f.png}
%%%END
%%%BLOCK id=224-08 ch=01 sec=匀变速直线运动推论 kind=formula alt=none cont=none y=0.86-0.92
平均速度\quad $V_{t/2}=\dfrac{V_0+V}{2}=V_0+\dfrac12at=\dfrac12(2V_0+at)=\dfrac12(V_0+V_0+at)$

中时速
%%%END
%%%BLOCK id=224-09 ch=01 sec=匀变速直线运动推论 kind=formula alt=none cont=none y=0.92-1.00
中位速\quad $V_{x/2}=\sqrt{\dfrac{V_0^2+V^2}{2}}$

%FIG p224_g 0.185 0.937 0.315 1.0
\notefig[0.25]{figures/p224_g.png}
$V_{x/2}>V_{t/2}$
\[
V'^2-V_0^2=2a\cdot\frac12x\qquad V^2-V'^2=2a\cdot\frac12x
\]
%FIG p224_h 0.34 0.947 0.545 1.0
\notefig[0.3]{figures/p224_h.png}
%%%END
%%%PAGEEND 224
